\chapter{What contributing means}
\label{ch:idea}

Open-source software is software whose code anyone may read and improve. Contributing means proposing a change to such a project. The change need not be code. Fixing a spelling mistake in the instructions counts, as does clearing up a confusing sentence, reporting a problem clearly, or adding a test. Beginners are often the best people to improve instructions, because they notice what the instructions assume.

You cannot damage a project by trying. Nothing you do reaches it until one of its maintainers, the people who look after it, reads your proposal and accepts it.

\section{The whole journey}

\begin{enumerate}
\item Find a project and a small thing to improve.
\item Make your own copy of the project on GitHub. This is a \emph{fork}.
\item Download your copy to your computer.
\item Make a \emph{branch}, which is a separate line of work.
\item Change the files, check them, and save a version.
\item Send the branch to your copy on GitHub.
\item Open a \emph{pull request}, asking the maintainers to take your change.
\item Answer their comments until the change is accepted or closed.
\end{enumerate}

\input{shared_git_words}

\input{shared_terminal}
